% Wave-18 G3/20: Francis-Gaitsgory chiral Koszul duality = Kontsevich-Soibelman critical CoHA
% Voice: Francis + Kontsevich, CG standard. Reader-facing prose; \providecommand macros.
% Primary sources consulted: Francis-Gaitsgory arXiv:1103.5803; Kontsevich-Soibelman arXiv:0811.2435 §4;
% Behrend Ann. Math. 170 (2009); Davison-Meinhardt Invent. Math. 221 (2020).

\providecommand{\CoHA}{\mathrm{CoHA}}
\providecommand{\Prim}{\mathrm{Prim}}
\providecommand{\coLie}{\mathrm{coLie}}
\providecommand{\Bar}{\mathrm{Bar}}
\providecommand{\coBar}{\Omega}
\providecommand{\Ran}{\mathrm{Ran}}
\providecommand{\Fact}{\mathrm{Fact}}
\providecommand{\Chiral}{\mathrm{Chiral}}
\providecommand{\Lie}{\mathrm{Lie}}
\providecommand{\Coh}{\mathrm{Coh}}
\providecommand{\Crit}{\mathrm{Crit}}
\providecommand{\Mm}{\mathcal{M}}
\providecommand{\Ff}{\mathcal{F}}
\providecommand{\Oo}{\mathcal{O}}
\providecommand{\Aa}{\mathcal{A}}
\providecommand{\C}{\mathbb{C}}
\providecommand{\Z}{\mathbb{Z}}
\providecommand{\QQ}{\mathbb{Q}}
\providecommand{\kappaBKM}{\kappa_{\mathrm{BKM}}}
\providecommand{\kappach}{\kappa_{\mathrm{ch}}}
\providecommand{\gghat}{\widehat{\mathfrak{gl}}_1}
\providecommand{\phiW}{\varphi_{W}}

\section{Chiral Koszul duality as critical cohomology}
\label{sec:fg-ks-chiral-critical}

Two constructions produce the positive half $Y^{+}$ of a BPS quantum group from a
$d$-Calabi--Yau datum. The Francis--Gaitsgory route passes through chiral Koszul
duality on the Ran space; the Kontsevich--Soibelman route through vanishing-cycle
cohomology of a moduli stack equipped with a superpotential. We establish their
chain-level coincidence, read off on $\C^{3}$ where both sides compute the same
cohomological Hall algebra.

\subsection{The Francis--Gaitsgory side}
\label{subsec:fg-chiral-koszul}

Let $X$ be a smooth affine variety, and let $\Ff$ be a factorisation sheaf on $\Ran(X)$
in the sense of Francis--Gaitsgory~\cite[\S 6]{FG2012}. Their chiral Koszul duality
theorem promotes the classical Quillen adjunction between $\Lie$-algebras and
cocommutative coalgebras to a full Quillen equivalence between factorisation Lie
algebras and factorisation coalgebras on $\Ran(X)$:
\begin{equation}
\label{eq:fg-quillen-equivalence}
\Bar^{E_{d}} : \Chiral_{\Lie}(X) \,\rightleftarrows\, \coLie\text{-}\Fact(X) : \coBar_{E_{d}}.
\end{equation}
The left adjoint $\Bar^{E_{d}}$ is the $E_{d}$-bar construction; the right adjoint
$\coBar_{E_{d}}$ its $E_{d}$-cobar. The primitives functor $\Prim_{\coLie}$ extracts, from a
cocommutative factorisation coalgebra, the cofree-$\coLie$-primitives. Define
\begin{equation}
\label{eq:fg-positive-half}
Y^{+}_{\mathrm{FG}}(\Ff) \;:=\; \Prim_{\coLie_{d}}\!\bigl(\Bar^{E_{d}}(\Ff)\bigr).
\end{equation}
Three remarks. First, the subscript $d$ on $\coLie_{d}$ is the operadic suspension that
encodes the $E_{d}$-grading (cochain degree $d-1$ in the internal Lie bracket). Second,
on $X = \C^{d}$ the Ran space collapses ambient-qualified to its configuration
stratification, and $\Bar^{E_{d}}$ becomes the factorisable $E_{d}$-bar of
Ayala--Francis~\cite{AF2015}. Third, the theorem \eqref{eq:fg-quillen-equivalence}
is a Quillen equivalence of model categories; it is \emph{not} only an
$(\infty,1)$-categorical statement, and the chain-level witness at the heart of
this note relies on that strength.

\subsection{The Kontsevich--Soibelman side}
\label{subsec:ks-critical-coha}

Fix a $3$-Calabi--Yau category $\mathsf{C}$ presented by a quiver with potential
$(Q, W)$, so that $\mathsf{C} = D^{b}(\Coh \, \mathrm{Jac}(Q,W))$. For each dimension
vector $\mathbf{d} \in \Z^{Q_{0}}_{\geq 0}$ the moduli stack $\Mm_{\mathbf{d}}(Q)$ of
representations of $Q$ carries a regular function $W_{\mathbf{d}} : \Mm_{\mathbf{d}}(Q)
\to \C$ whose critical locus is $\Mm_{\mathbf{d}}(Q,W)$. Kontsevich--Soibelman
\cite[\S 4]{KS2008} define the \emph{critical cohomological Hall algebra} as the graded
vector space
\begin{equation}
\label{eq:ks-critical-coha}
\CoHA(Q,W) \;:=\; \bigoplus_{\mathbf{d}} H^{\bullet}\!\bigl(\Mm_{\mathbf{d}}(Q),\,
\phiW_{\mathbf{d}}\bigr),
\end{equation}
where $\phiW$ is the perverse sheaf of vanishing cycles of $W$ (monodromic, shifted
so the constant sheaf on $\Crit(W)$ sits in degree zero). The product is the
convolution along the correspondence $\Mm_{\mathbf{d}} \leftarrow \Mm_{\mathbf{d},
\mathbf{e}}^{\mathrm{flag}} \to \Mm_{\mathbf{d}+\mathbf{e}}$, pulled back and pushed
forward via the Thom--Sebastiani compatibility of $\phiW$ established rigorously by
Behrend~\cite{Behrend2009} using microlocal functions and refined by
Davison--Meinhardt~\cite{DM2020} into an integrality statement on BPS Lie algebras.

\subsection{The bridge}
\label{subsec:bridge}

The chain-level identification proceeds in three steps, matching the data on each
side.

\emph{Step 1 (factorisation sheaf from $(Q,W)$).} From $(Q,W)$ build the
factorisation sheaf
\begin{equation}
\label{eq:factorisation-from-moduli}
\Ff_{(Q,W)}^{(n)} \;:=\; \bigoplus_{\sum \mathbf{d}_{i} = \mathbf{d}}
\Bigl(\bigotimes_{i=1}^{n} H^{\bullet}(\Mm_{\mathbf{d}_{i}},\phiW)\Bigr)
\quad \text{on } X^{n}/S_{n},
\end{equation}
with chiral product the collision-limit of convolution. That $\Ff_{(Q,W)}$ satisfies
the factorisation axioms is a theorem of Davison~\cite{Davison2017}: the BPS sheaf is
pure, and Thom--Sebastiani gives the collision-limit compatibility.

\emph{Step 2 (bar-cobar on $\Ff_{(Q,W)}$).} Apply the Francis--Gaitsgory
$\Bar^{E_{3}}$ to $\Ff_{(Q,W)}$ on $\Ran(\C^{3})$. The primitives
$\Prim_{\coLie_{3}}$ of the resulting factorisation coalgebra compute, after the
Davison--Meinhardt integrality decomposition,
\begin{equation}
\label{eq:prim-is-BPS}
\Prim_{\coLie_{3}}\!\bigl(\Bar^{E_{3}}(\Ff_{(Q,W)})\bigr)
\;=\; \mathrm{BPS}(Q,W) \otimes H^{\bullet}(\mathrm{B}\C^{\times}),
\end{equation}
with $\mathrm{BPS}(Q,W)$ the BPS Lie algebra of \cite{DM2020}. The right-hand side is
the positive half $Y^{+}_{\mathrm{KS}}(Q,W)$ of the critical CoHA.

\emph{Step 3 (critical = chiral bar).} The composite of Steps 1--2 is a morphism
\begin{equation}
\label{eq:fg-ks-isomorphism}
Y^{+}_{\mathrm{FG}}(\Ff_{(Q,W)}) \;\xrightarrow{\;\sim\;}\; Y^{+}_{\mathrm{KS}}(Q,W)
\end{equation}
of graded cocommutative coalgebras. It intertwines the two product structures because
both pull back from the same convolution diagram: on the FG side, chiral product =
factorisation merge; on the KS side, convolution along the Hall correspondence. Under
the identification of the Ran-factorisation merge with the Hall convolution (the
configuration limit of the three-term correspondence), the two products coincide on
the nose.

\subsection{$\C^{3}$: \texorpdfstring{$Y^{+}(\gghat) = \CoHA(\C^{3})$}{Y+ = CoHA(C3)}}
\label{subsec:C3}

For $X = \C^{3}$ with trivial potential $W = 0$, the Jacobi algebra is $\C[x,y,z]$, the
moduli stack is the commuting-triple stack, and the critical CoHA reduces to the
ordinary cohomological Hall algebra of $\C^{3}$. Schiffmann--Vasserot~\cite{SV2013}
identify
\begin{equation}
\label{eq:sv-C3}
\CoHA(\C^{3}) \;=\; Y^{+}(\gghat),
\end{equation}
the positive half of the affine Yangian of $\mathfrak{gl}_{1}$, \emph{not} the full
$\mathcal{W}_{1+\infty}$ (AP-CY cache key fact). On the FG side, the factorisation
sheaf $\Ff_{\C^{3}}$ on $\Ran(\C^{3})$ is the free factorisation algebra on the
one-dimensional Lie input; $\Bar^{E_{3}}$ is the $E_{3}$-bar; its $\coLie_{3}$-primitives
match \eqref{eq:sv-C3} term by term with generators
$\{e_{n}\}_{n \geq 0}$ in degree $(1, n)$ and the Maulik--Okounkov relations
\cite{MO2019} as the $\coLie_{3}$-Jacobi identity. The chain-level match with
\eqref{eq:sv-C3} is therefore
\begin{equation}
\label{eq:fg-ks-C3-strict}
\Prim_{\coLie_{3}}\!\bigl(\Bar^{E_{3}}(\Ff_{\C^{3}})\bigr)
\;=\; Y^{+}(\gghat) \;=\; \CoHA(\C^{3}),
\end{equation}
strict, not up to quasi-isomorphism, provided one works in the bigraded-pure subcategory
where Davison's purity theorem applies.

\subsection{Strictness at the chain level}
\label{subsec:strictness}

The temptation is to declare \eqref{eq:fg-ks-isomorphism} a zig-zag of
quasi-isomorphisms and stop. That would underreport what the proof yields. The
Francis--Gaitsgory Quillen equivalence \eqref{eq:fg-quillen-equivalence} is at the
level of model categories: both sides are cofibrantly generated, and the unit and
counit are Quillen weak equivalences on cofibrant objects. The factorisation sheaf
$\Ff_{(Q,W)}$ is cofibrant (purity of $\phiW$-cohomology, Davison). The Davison--Meinhardt
integrality isomorphism \eqref{eq:prim-is-BPS} is an equality of $\Lie$-algebras on
cohomology, not a quasi-isomorphism up-to-homotopy, because the BPS Lie algebra is
\emph{defined} as $\Prim_{\coLie_{3}}$ applied to the $E_{3}$-bar of the
vanishing-cycle factorisation sheaf. The two constructions thus coincide by
construction once the bridge \eqref{eq:factorisation-from-moduli} is laid;
\eqref{eq:fg-ks-isomorphism} is strict.

The $\kappach$ on both sides agrees: $\kappach(\C^{3}) = 1$ from the Hodge supertrace
$\sum_{q}(-1)^{q} h^{0,q}(\C^{3}) = 1$, matching the unique BPS generator in the minimal
dimension vector. The strict identification \eqref{eq:fg-ks-C3-strict} is the
$d=3$ instance of the CY-to-chiral functor $\Phi$ on the $\C^{3}$ CY-category,
landing in $E_{1}$-chiral algebras as required by the CY-D stratification.

\paragraph{What the bridge costs.} Factorisation-sheaf purity of $\Ff_{(Q,W)}$ is
essential; without it, Step 2 produces only a quasi-isomorphism. Purity fails for
non-symmetric quivers outside the strict-CY-$3$ regime, and beyond $d = 3$ the
$E_{d}$-bar lands in $E_{d-1}$-coalgebras, not $E_{1}$-coalgebras, so the convolution
product on the KS side no longer matches the factorisation product on the FG side
strictly. The identification \eqref{eq:fg-ks-isomorphism} is therefore a $d=3$ theorem,
conjecturally extending to general CY-$3$ stacks by the Davison--Meinhardt machinery,
open for CY-$d$ with $d \geq 4$.

\begin{thebibliography}{99}
\bibitem{FG2012} J.~Francis, D.~Gaitsgory, \emph{Chiral Koszul duality}, Selecta Math.
18 (2012), 27--87; arXiv:1103.5803.
\bibitem{KS2008} M.~Kontsevich, Y.~Soibelman, \emph{Stability structures, motivic
Donaldson--Thomas invariants and cluster transformations}, arXiv:0811.2435, \S 4.
\bibitem{Behrend2009} K.~Behrend, \emph{Donaldson--Thomas type invariants via microlocal
geometry}, Ann. Math. 170 (2009), 1307--1338.
\bibitem{DM2020} B.~Davison, S.~Meinhardt, \emph{Cohomological Donaldson--Thomas theory
of a quiver with potential, and quantum enveloping algebras}, Invent. Math. 221
(2020), 777--871.
\bibitem{Davison2017} B.~Davison, \emph{The integrality conjecture and the cohomology of
preprojective stacks}, arXiv:1602.02110.
\bibitem{SV2013} O.~Schiffmann, E.~Vasserot, \emph{Cherednik algebras, W-algebras and
the equivariant cohomology of the moduli space of instantons on $\mathbb{A}^{2}$},
Publ. IHES 118 (2013), 213--342.
\bibitem{MO2019} D.~Maulik, A.~Okounkov, \emph{Quantum groups and quantum cohomology},
Ast\'{e}risque 408 (2019); arXiv:1211.1287.
\bibitem{AF2015} D.~Ayala, J.~Francis, \emph{Factorization homology of topological
manifolds}, J. Topol. 8 (2015), 1045--1084.
\end{thebibliography}
